\documentclass[10pt]{article}
\usepackage[margin=0.85in]{geometry}
\usepackage{amsmath,amssymb,amsthm}
\usepackage{bm}
\usepackage{booktabs}
\usepackage[colorlinks=true,linkcolor=blue,citecolor=blue,urlcolor=blue]{hyperref}

\newcommand{\R}{\mathbb{R}}
\newcommand{\pdd}[2]{\frac{\partial #1}{\partial #2}}
\newcommand{\pddt}[2]{\frac{\partial^{2} #1}{\partial #2^{2}}}
\newcommand{\erfc}{\operatorname{erfc}}
\newcommand{\norm}[1]{\left\lVert #1 \right\rVert}
\newcommand{\Cvar}{\mathrm{CVaR}}
\newcommand{\E}{\mathbb{E}}

\title{\vspace{-1.2em}Equations and Algorithms Under Test:\\ Groundwater Remediation via PINNs, FNOs, and Stochastic Optimal Control}
\author{S. Lucero}
\date{\today}

\begin{document}
\maketitle
\vspace{-2.5em}

\begin{abstract}
\noindent This note records the mathematical models and numerical algorithms
currently exercised in the repository. Two families are active: (i) the
physics-informed neural network (PINN) and Fourier neural operator (FNO)
surrogates for the advection--dispersion equation, and (ii) the
finite-volume / sample-average-approximation (SAA) stochastic optimal control
pipeline for pump-and-treat remediation. We state the governing equations, the
learning and optimization objectives, and the solver algorithms exactly as
they are implemented, so the manuscript can be rebuilt from this single source.
\end{abstract}

\section{Governing equations}
\label{sec:pde}

\paragraph{1D advection--dispersion (sandbox).}
The dimensionless, homogeneous, non-reactive 1D problem used to verify the
neural surrogates is
\begin{equation}
\pdd{u}{\tau}+\pdd{u}{\xi}=\frac{1}{\mathrm{Pe}}\,\pddt{u}{\xi},
\quad (\xi,\tau)\in(0,1)\times(0,\tau_{\max}],
\label{eq:ade1d}
\end{equation}
with $u(\xi,0)=0$, $u(0,\tau)=1$, $u(1,\tau)=0$, where $\xi=x/L$,
$\tau=tv/L$ (pore volumes), $u=C/C_0$, and $\mathrm{Pe}=vL/D$. The
closed-form Ogata--Banks reference \cite{ogata1961solution},
\begin{equation}
u_{\mathrm{ref}}(\xi,\tau)=
\frac12\Bigl[
\erfc\!\Bigl(\tfrac{(\xi-\tau)\sqrt{\mathrm{Pe}}}{2\sqrt{\tau}}\Bigr)
+e^{\mathrm{Pe}\,\xi}\erfc\!\Bigl(\tfrac{(\xi+\tau)\sqrt{\mathrm{Pe}}}{2\sqrt{\tau}}\Bigr)
\Bigr],
\label{eq:ogata}
\end{equation}
is the ground truth for the surrogate tests.

\paragraph{2D flow (steady Darcy).}
The pump-and-treat physics is depth-averaged over a rectangular domain
$\Omega\subset\R^2$ with cell-centred grid. The steady flow equation is
\begin{equation}
-\nabla\cdot\bigl(T(\mathbf{x})\,\nabla h\bigr)
= b\,W(\mathbf{x})-\sum_{i=1}^{n_w} q_i\,\delta_\varepsilon(\mathbf{x}-\mathbf{x}_i),
\label{eq:darcy}
\end{equation}
with transmissivity $T=K b$ ($K$ hydraulic conductivity, $b$ thickness), head
$h$~[m], recharge $W$, and extraction rates $q_i\ge0$~[$\mathrm{m^3\,d^{-1}}$].
Wells are regularized by a normalized Gaussian kernel of width $\varepsilon$
grid cells,
\begin{equation}
\delta_\varepsilon(\mathbf{x}-\mathbf{x}_i)=\frac1{Z_i}
\exp\!\Bigl(-\tfrac{(x-x_i)^2}{2\sigma_x^2}-\tfrac{(y-y_i)^2}{2\sigma_y^2}\Bigr),
\qquad \sigma_\bullet=\varepsilon\,\Delta_\bullet .
\label{eq:well}
\end{equation}
Boundary conditions are Dirichlet heads on the $x$-faces and no-flow (or
Dirichlet) on the $y$-faces. The seepage velocity is
$\mathbf{v}=-K\,\nabla h/\phi$ with porosity $\phi$.

\paragraph{2D transport (implicit Euler).}
Concentration obeys the depth-averaged advection--dispersion equation with a
frozen (background) dispersion tensor and explicit well sinks,
\begin{equation}
\phi R\,\pdd{C}{t}
=\nabla\cdot(\phi\,\mathbf{D}\,\nabla C)
-\nabla\cdot(\phi\,\mathbf{v}\,C)
-\sum_i q_i\,C(\mathbf{x}_i)\,\delta_\varepsilon(\mathbf{x}-\mathbf{x}_i)
+s(\mathbf{x},t),
\label{eq:transport}
\end{equation}
with retardation $R$ and source $s$. Time is discretized by implicit Euler,
and advection uses first-order upwinding, so the transport operator is linear
in $q$ and amenable to an exact discrete adjoint.

\section{Neural surrogates}
\label{sec:surrogates}

\paragraph{PINN.}
The solution of \eqref{eq:ade1d} is approximated by a multilayer perceptron
$u_\theta(\xi,\tau)$ ($\tanh$ hidden, sigmoid output). The PDE residual
$\mathcal{R}_\theta=\pdd{u_\theta}{\tau}+\pdd{u_\theta}{\xi}-\tfrac{1}{\mathrm{Pe}}\pddt{u_\theta}{\xi}$
is formed by automatic differentiation, and $\theta$ minimizes the composite
loss
\begin{equation}
\mathcal{L}(\theta)=\lambda_{\mathrm{PDE}}\mathcal{L}_{\mathrm{PDE}}
+\lambda_{\mathrm{IC}}\mathcal{L}_{\mathrm{IC}}
+\lambda_{\mathrm{BC}}\mathcal{L}_{\mathrm{BC}},
\label{eq:pinnloss}
\end{equation}
where each term is a mean-squared residual over randomly resampled collocation
points. A \emph{hard} (architectural) ansatz enforces IC and outlet exactly:
\begin{equation}
u_\theta(\xi,\tau)=(1-\xi)\,s(\tau)+\xi(1-\xi)\,s(\tau)\,\mathrm{NN}_\theta(\xi,\tau),
\quad s(\tau)=1-e^{-\tau/\varepsilon},
\label{eq:hardbc}
\end{equation}
with $\varepsilon=10^{-3}$. Optimization is Adam ($\sim\!10^4$ steps with
collocation resampling) followed by L-BFGS. A known artefact: the inlet ramp
contributes $s'(0)=1/\varepsilon$, so $\max|\mathcal{R}_\theta|\gtrsim 1/\varepsilon$
independently of the weights; the entrance corner $(0,0)$ is likewise a
measure-zero discontinuity excluded from the reported metric.

\paragraph{FNO.}
The FNO learns the solution operator $\mathcal{G}:a\mapsto u$ for a family of
hydrologies $a=(v,D)$ or $(\mathrm{Pe})$, discretized as
$u_j^n\approx\mathcal{G}_\theta(a;\xi_j,\tau_n)$. It is a lifting MLP, $K$
spectral layers with $k_{\max}$ retained Fourier modes, and a projection MLP
\cite{li2020fourier}. A structural requirement discovered in Phase~4 is that
normalized $(\xi,\tau)$ coordinate channels must be appended to the input:
with only spatially constant channels the input has a single (DC) mode and the
operator output is constant by construction (measured profile-only error
$1.0$). With coordinate channels the test relative $L^2$ error reaches
$\sim\!1.5\times10^{-2}$ at 250 epochs.

\section{Stochastic optimal control problem}
\label{sec:saa}

Heterogeneity is represented by $M$ random conductivity fields
$\{K_m\}_{m=1}^M$ (Monte Carlo scenarios). For a policy
$q\in[0,q_{\max}]^{n_w}$ with plant capacity $Q_{\max}$, the per-scenario cost
is
\begin{equation}
J_m(q)=J_{\mathrm{pump}}+J_{\mathrm{treat}}+J_{\mathrm{res}}+J_{\mathrm{time}},
\label{eq:cost}
\end{equation}
with
\begin{align}
J_{\mathrm{pump}}&=c_e\,T\sum_i q_i,
&
J_{\mathrm{treat}}&=c_w\sum_i q_i\!\int_0^T C_{\mathrm{well},i}(t)\,dt,\notag\\[2pt]
J_{\mathrm{res}}&=c_m\!\int_\Omega b R\phi\,C(\mathbf{x},T)\,d\Omega,
&
J_{\mathrm{time}}&=c_\tau\,\tau_c+c_{\mathrm{fail}}\,\mathbf{1}_{\{\text{failed}\}}.
\label{eq:costterms}
\end{align}
Here $\tau_c$ is the cleanup time (last time the compliance violation is
positive), replaced by a differentiable sigmoid surrogate when gradients are
needed. Two chance constraints bind the policy: a compliance limit at the
monitoring well,
\begin{equation}
G_m(q)=\max_{t\in[t_0,T]}\bigl(C_{\mathrm{comp}}(t)-C_{\mathrm{reg}}\bigr)\le0
\quad\text{w.p.}\ge1-\alpha,
\label{eq:G}
\end{equation}
and a drawdown limit at observation wells,
$H_m(q)=\min(h_{\mathrm{obs}}-h_{\min})\ge0$ w.p. $\ge1-\alpha_h$. The
nonsmooth empirical chance constraint is replaced by the conservative CVaR
reformulation \cite{rockafellar2000optimization},
\begin{equation}
\Cvar_{1-\alpha}(G)=\min_{\eta}\Bigl\{\eta+\tfrac{1}{\alpha M}\sum_m [G_m-\eta]_+\Bigr\}\le0,
\label{eq:cvar}
\end{equation}
and the penalized SAA objective is
\begin{equation}
f(q)=\frac1M\sum_m J_m(q)
+\rho_g\,[\Cvar_{1-\alpha}(G)]_+^2
+\rho_h\,[\Cvar_{1-\alpha_h}(-H)]_+^2,
\label{eq:saa}
\end{equation}
with penalty weights $\rho_g,\rho_h$ optionally continued outward.

\section{Numerical algorithms}
\label{sec:alg}

\paragraph{Finite-volume discretization.}
Flow \eqref{eq:darcy} is discretized on a cell-centred five-point stencil with
harmonic-mean face transmissivity, giving a symmetric positive-definite (SPD)
system $\mathbf{A}\mathbf{h}=\mathbf{b}(q)$. The transport operator in
\eqref{eq:transport} is assembled as an upwind advection stencil plus a
frozen-dispersion stencil; the well sink uses the kernel \eqref{eq:well}.
Implicit Euler gives, per time step,
\begin{equation}
\Bigl(\tfrac{\phi R}{\Delta t}\mathbf{I}+\mathbf{A}_{\mathrm{adv}}(q)+\mathbf{A}_{\mathrm{disp}}\Bigr)
\mathbf{C}^{n+1}=\tfrac{\phi R}{\Delta t}\mathbf{C}^{n}+\mathbf{s}^{n}.
\label{eq:ie}
\end{equation}

\paragraph{Batched Krylov solvers (PyTorch backend).}
The NumPy/SciPy reference uses a sparse LU factorization per scenario per step.
The GPU backend keeps the identical algebra but changes the execution model to
matrix-free stencils and batched iterative solves over $S\times M$ systems
($S$ candidate policies, $M$ scenarios): conjugate gradients (CG) for the SPD
Darcy system, and BiCGSTAB for the nonsymmetric transport system, both with
Jacobi preconditioning. The transport loop is warm-started from the previous
time level, giving $8$--$9$ BiCGSTAB iterations per step at $rtol=10^{-11}$
(float64) versus $\sim\!17$ cold.

\paragraph{Optimization.}
The outer loop minimizes \eqref{eq:saa} with capacity projection
$\mathrm{proj}(q)=\min(q,q_{\max})$, scaled by $Q_{\max}$ if
$\sum_i q_i>Q_{\max}$. Two deterministic search strategies coexist:
(i) \emph{compass pattern search}, which probes $q\pm\delta e_i$ one candidate
at a time and halves the step on failure; and (ii) \emph{batched pattern
search}, which evaluates the whole compass set
$\{q\pm\delta e_i\}$ in one batched call and accepts the best strict
improvement --- on a GPU the marginal cost of the extra candidates is near zero.
For gradient-based search, SLSQP uses a discrete-adjoint gradient of the
smooth (CVaR + soft-cleanup) objective, exploiting that the frozen-dispersion
upwind operator is linear in $q$. Multi-start ($n_{\mathrm{starts}}$ random
feasible seeds) keeps the best penalized objective. The measured bounded run
($80\times40$ grid, $M=40$ train / $80$ test, $200$ evaluations) terminates in
$\sim\!19$ min at $\approx5$ s per evaluation, converging in $\sim\!150$
evaluations.

\begin{thebibliography}{9}
\bibitem{ogata1961solution}
A. Ogata and R. B. Banks, \textit{A solution of the differential equation of
longitudinal dispersion in porous media}, U.S. Geological Survey Professional
Paper 411-A, 1961.
\bibitem{li2020fourier}
Z. Li, N. Kovachki, K. Azizzadenesheli, B. Liu, K. Bhattacharya, A. Stuart,
A. Anandkumar, \textit{Fourier Neural Operator for Parametric Partial
Differential Equations}, ICLR, 2021.
\bibitem{rockafellar2000optimization}
R. T. Rockafellar and S. Uryasev, \textit{Optimization of conditional
value-at-risk}, Journal of Risk, 2:21--41, 2000.
\end{thebibliography}

\end{document}
